% ─────────────────────────────────────────────────────────────────────────────
\section{EN 1990~: les combinaisons d'actions}
% ─────────────────────────────────────────────────────────────────────────────

\begin{maitre}
Avant de v\'erifier une poutre, il faut savoir sous quelles charges
on la v\'erifie. Qu'est-ce qui charge une panne de toiture~?
\end{maitre}

\begin{apprenti}
Le poids propre du bois, le poids de la couverture, la neige\ldots{}
et peut-\^etre le vent~?
\end{apprenti}

\begin{maitre}
Exactement. Ces charges s'appellent des \textbf{actions} en langage Eurocode.
Et la norme EN~1990 dit qu'on ne peut pas simplement les additionner~---
il faut construire des \textbf{combinaisons} avec des coefficients
de s\'ecurit\'e, parce que la probabilit\'e que toutes les charges
soient simultan\'ement \`a leur valeur maximale est tr\`es faible.
\end{maitre}

La formule g\'en\'erale de combinaison ELU (\'Etat Limite Ultime ---
pour la r\'esistance) selon l'AN France est~:

\begin{equation}
E_{d} = \underbrace{\gamma_G \cdot G_k}_{\text{permanent}}
      + \underbrace{\gamma_{Q,1} \cdot Q_{k,1}}_{\text{action principale}}
      + \underbrace{\sum_{i>1} \psi_{0,i} \cdot \gamma_{Q,i} \cdot Q_{k,i}}_{\text{actions d'accompagnement}}
\label{eq:elu}
\end{equation}

\noindent o\`u~:
\begin{itemize}
  \item $G_k$ = charge permanente caract\'eristique (kN/m\textsuperscript{2})
  \item $Q_{k,1}$ = action variable principale (neige, vent ou exploitation)
  \item $Q_{k,i}$ = actions variables d'accompagnement
  \item $\gamma_G = 1{,}35$ (charges permanentes d\'efavorables)
  \item $\gamma_Q = 1{,}50$ (charges variables)
  \item $\psi_0$ = coefficient de combinaison (r\'eduit les actions secondaires)
\end{itemize}

Le code g\'en\`ere plusieurs cas~: une fois avec la neige dominante
(\py{ELU\_S}), une fois avec le vent dominant (\py{ELU\_W}),
une fois charges permanentes seules (\py{ELU\_G}), etc.

\begin{lstlisting}[style=python, caption={Extrait de \fichier{ec0/combinaisons.py} --- coefficients AN France et génération}]
# Coefficients réglementaires AN France
GAMMA_G = 1.35    # charges permanentes défavorables
GAMMA_Q = 1.50    # charges variables

# Coefficients ψ₀ : réduction des actions d'accompagnement
PSI_0_S = 0.5     # neige (altitude <= 1000 m)
PSI_0_W = 0.6     # vent
PSI_0_Q = 0.7     # exploitation catégorie A/B (habitation)

# Durées de charge (pour le facteur k_mod EC5 Tableau 3.1)
# G      -> permanent
# neige  -> court_terme
# vent   -> instantane

# Génération de la combinaison ELU avec la neige comme action principale
if s_k > 0:   # seulement s'il y a une charge de neige configurée
    combinaisons.append(CombinaisonEC0(
        id_combinaison  = "ELU_S",
        type_combinaison= "ELU_STR",
        charge_principale="S",
        gamma_G         = GAMMA_G,     # 1.35
        gamma_Q1        = GAMMA_Q,     # 1.50 * neige
        psi_0_Q2        = PSI_0_W if w_k > 0 else 0.0,  # vent accompagne
        duree_charge    = "court_terme",
    ))

# Et de même pour ELU_W (vent principal), ELU_Q (exploitation principale)...
\end{lstlisting}

Pour les \textbf{ELS} (\'Etats Limites de Service --- pour les fl\`eches),
les m\^emes actions sont combin\'ees \emph{sans} coefficients partiels~:

\begin{equation}
E_{d,ELS,car} = G_k + Q_{k,1} + \sum_{i>1} \psi_{0,i} \cdot Q_{k,i}
\label{eq:els_car}
\end{equation}